\chapter{mmap、page cache 与文件映射}

\section{原理}

\code{read} 和 \code{write} 把文件内容显式复制到用户缓冲区，而 \code{mmap} 把文件区间映射到进程地址空间，让普通 load/store 触发文件页面的调入和回写。它看起来像内存功能，实际是虚拟内存、文件系统和 page cache 的共同契约：VMA 记录虚拟地址区间；page cache 记录文件偏移到物理页的缓存；缺页处理把两者连接起来；写回系统决定脏页何时落盘。

先写一个常见误解：

\begin{lstlisting}
addr = mmap(file, length)
# wrong mental model:
# kernel reads the whole file into memory now
# addr points to that private copy
\end{lstlisting}

这个模型会立刻解释不了三个现象。第一，映射一个很大的文件可以很快返回，因为内核通常只建立 VMA，不会立刻把所有页面读入内存。第二，多个进程映射同一个文件区间，读到的可能是同一份 page cache 页，而不是每个进程一份私有拷贝。第三，另一个进程把文件 truncate 后，当前进程访问原来映射的越界部分可能得到 \code{SIGBUS}，因为虚拟地址仍在 VMA 里，但背后的文件页已经不存在。

正确过渡是：

\begin{lstlisting}
mmap:
  create VMA [addr, addr + length)
  remember file and file offset
  do not install every PTE yet

first load from addr + x:
  page fault
  find VMA
  compute file page index
  find or load page cache page
  install PTE
resume user instruction
\end{lstlisting}

所以 \code{mmap} 的核心不是“少一次拷贝”这一句话，而是把文件偏移延迟绑定到虚拟地址访问上。这个延迟绑定带来性能，也带来错误边界：缺页可能等待 I/O，写入可能变成 COW 或 dirty page，持久化要靠 \code{msync/fsync} 推进，文件大小变化可能让合法 VMA 里的访问失败。

用一个小文件把状态走完整。文件 \code{data.bin} 有 9000 字节，页大小 4096。进程执行：

\begin{lstlisting}
fd = open("data.bin", O_RDWR);
p = mmap(NULL, 8192, PROT_READ | PROT_WRITE,
         MAP_SHARED, fd, 0);
x = p[5000];
\end{lstlisting}

\code{mmap} 返回时，常见状态并不是“两个页都已经在内存里”。更准确的账本是：

\begin{longtable}{p{0.18\linewidth}p{0.34\linewidth}p{0.34\linewidth}}
\toprule
对象 & 刚 mmap 后 & 第一次读 \code{p[5000]} 后 \\
\midrule
VMA & 记录虚拟区间、文件、offset 0、共享可写 & 不变 \\
PTE & 两个虚拟页通常都还没有 present 映射 & 第二个虚拟页 present，指向 page cache 页 \\
page cache & 可能没有对应页，也可能已有别的进程留下的页 & \code{(inode, index=1)} 有缓存页 \\
文件系统 I/O & 通常没有立即读 8192 字节 & 若 cache miss，提交读第 1 个文件页 \\
用户指令 & 还没执行 load & fault 返回后重新执行 load，读到字节 \\
\bottomrule
\end{longtable}

这里的 \code{p[5000]} 落在映射的第二个虚拟页，因为 \code{5000 / 4096 = 1}。缺页处理先用 fault 地址找到 VMA，再把 \code{addr - vma.start} 换算成文件页索引 1，然后到 inode 的 page cache 中查 \code{index=1}。若缓存命中，内核直接安装 PTE；若缓存未命中，内核分配一个 locked page，把它先放入 page cache，再发起磁盘读。读完成前 faulting 线程睡眠；完成后解锁页面，安装 PTE，回到用户态重新执行那条 load。注意是重新执行 faulting 指令，不是从下一行继续，因为第一次 load 还没有真正完成。

再看边界。映射长度是 8192，只覆盖文件前两个页；文件实际有 9000 字节，所以这两个页都在文件范围内。若改成映射 16384 字节，再访问第三页的一部分，VMA 可能仍然覆盖该地址，但文件没有对应内容，内核不能凭空制造稳定文件页，应该按系统语义触发 \code{SIGBUS} 或等价错误。这也是 \code{SIGSEGV} 和 \code{SIGBUS} 的区别：前者常表示地址不属于合法映射或权限不对；后者常表示地址形式上属于映射，但背后的文件对象不能提供该页。

为什么需要 \code{mmap}？第一，它减少系统调用和复制，适合随机访问大文件。第二，它让多个进程共享同一份文件页，天然支持共享库和共享内存。第三，它把文件 I/O 纳入页表权限和缺页机制，统一了“按需加载”。但它也带来更复杂的一致性：一个进程通过 \code{write} 修改文件，另一个进程通过 \code{mmap} 读同一区间，看到的应是同一份 page cache 语义，而不是两个互相独立的缓存。

\begin{keyidea}
\code{mmap} 不是把文件一次性读进内存，而是建立“虚拟地址 -> 文件偏移 -> page cache 页面”的延迟绑定。
\end{keyidea}

\code{MAP\_SHARED} 和 \code{MAP\_PRIVATE} 是理解文件映射的分水岭。共享映射的写入会把 page cache 页面标脏，之后由 \code{msync}、\code{fsync} 或后台写回落盘；私有映射初始可共享文件页，但第一次写入触发 COW，产生匿名页，文件本身不被修改。共享库代码段通常是只读文件映射；进程堆和栈通常是匿名映射；数据库会谨慎选择 \code{mmap} 或 direct I/O，因为错误的写回假设会破坏延迟和持久化边界。

\section{实践}

一个教学 \code{mmap} 需要以下对象：

\begin{longtable}{p{0.2\linewidth}p{0.34\linewidth}p{0.34\linewidth}}
\toprule
对象 & 保存什么 & 不变量 \\
\midrule
VMA & 虚拟区间、权限、flags、file、offset & VMA 不重叠，权限不超过文件打开模式 \\
page table entry & 虚拟页到物理页、权限、dirty/accessed & PTE 权限不能超过 VMA 权限 \\
page cache page & inode、文件页索引、物理页、dirty、refcount & 同一 inode+index 对应一个缓存页 \\
anon page & 私有 COW 后的页面 & 不再回写到文件 \\
writeback item & 脏页、目标块、提交状态 & 写回完成前页面不能被当成干净页回收 \\
\bottomrule
\end{longtable}

文件映射缺页路径如下：

\begin{lstlisting}
handle_page_fault(addr, access):
  vma = find_vma(current.mm, addr)
  if vma == null or access not allowed by vma:
    send_sigsegv()
  page_index = vma.file_offset_pages + ((addr - vma.start) / PAGE_SIZE)
  page = page_cache_get_or_load(vma.file, page_index)
  if access == WRITE and vma.private:
    page = copy_page(page)
    mark_anon(page)
  install_pte(addr, page, pte_prot(vma, access))
\end{lstlisting}

\code{mmap} 实现中最重要的错误路径是长度、偏移和权限检查。映射长度为零、地址未对齐、偏移未按页对齐、写共享映射但 fd 没有写权限、文件系统不支持映射、区间越过用户地址空间，都必须明确返回错误。不要把所有失败都压成一个 \code{EINVAL}，否则测试和调试会失去边界信息。

\section{算法与实现分析}

\subsection{VMA 查找}

最简单的 VMA 管理可以用按起始地址排序的数组或链表，查找复杂度为 \code{O(n)}。这对教学模型足够，但进程有大量映射时会拖慢 page fault。真实内核会使用树结构，Linux 近年使用 maple tree 管理 VMA，以降低查找、插入和遍历成本。

\begin{lstlisting}
find_vma(mm, addr):
  for vma in mm.vmas_by_start:
    if vma.start <= addr and addr < vma.end:
      return vma
    if addr < vma.start:
      return null
  return null
\end{lstlisting}

插入 VMA 时必须检查相邻区间是否可合并。若权限、flags、file 和连续 offset 一致，就可以合并，减少 VMA 数量。若频繁切分和合并出错，会导致 \code{munmap} 后仍然能访问、或合法区间被错误 SIGSEGV。

\subsection{page cache 索引}

page cache 的键通常是 \code{(inode, page\_index)}。这个键保证 \code{read/write} 与 \code{mmap} 共享同一页：

\begin{lstlisting}
page_cache_get_or_load(file, index):
  key = (file.inode, index)
  lock(page_cache.lock)
  page = radix_tree_lookup(page_cache, key)
  if page != null:
    page.refcount++
    unlock(page_cache.lock)
    return page
  page = allocate_page_locked()
  radix_tree_insert(page_cache, key, page)
  unlock(page_cache.lock)
  read_file_page(file, index, page)
  unlock_page(page)
  return page
\end{lstlisting}

这里先插入 locked page，再执行 I/O，是为了阻止多个线程同时读取同一文件页。后来的缺页者发现页面存在但 locked，就等待页面解锁。否则并发 page fault 会造成重复 I/O，甚至最后一个完成者覆盖前一个完成者的状态。

\subsection{shared 与 private 写故障}

写故障要区分三种情况：VMA 不允许写，私有映射需要 COW，共享映射需要让 PTE 可写并跟踪脏页。

\begin{lstlisting}
handle_write_fault(vma, addr, old_page):
  if not vma.writable:
    send_sigsegv()
  if vma.private:
    new_page = allocate_page()
    copy(new_page, old_page)
    install_pte(addr, new_page, WRITE | USER)
    put_page(old_page)
    return
  mark_page_dirty(old_page)
  install_pte(addr, old_page, WRITE | USER | SHARED)
\end{lstlisting}

COW 的复杂度是一次页面复制 \code{O(PAGE\_SIZE)}。它用写时成本换取 fork/exec 和私有映射的低启动成本。共享映射的复杂性则在持久化：PTE dirty、page dirty、文件系统 journal dirty 和块设备队列不是同一件事。用户调用 \code{msync(MS\_SYNC)} 期望脏页写回并等待完成，但这仍不必然等价于包含目录项的业务提交，后者需要文件系统和上层协议共同保证。

继续用 \code{data.bin}。两个进程 A 和 B 都用 \code{MAP\_SHARED} 映射文件第一页。A 执行：

\begin{lstlisting}
p[10] = 'X';
\end{lstlisting}

如果 A 的 PTE 当前是只读但 VMA 允许写，CPU 会产生写权限 fault。内核确认这是共享可写映射后，可以把同一份 page cache 页装成可写 PTE，并在合适的时机把页面标脏。此后 B 通过 \code{read(fd, buf, ...)} 或自己的共享映射读同一文件页，看到的可能就是 page cache 中已修改的内容。此时状态应分层看：

\begin{longtable}{p{0.22\linewidth}p{0.58\linewidth}}
\toprule
层次 & \code{p[10]='X'} 之后发生了什么 \\
\midrule
CPU/PTE & A 的页表项允许写；硬件 dirty 位可能被置位 \\
page cache & \code{(inode,0)} 对应页内容已变，页面需要写回 \\
文件系统 & 还未必分配或提交所有元数据事务 \\
块层/设备 & 还未必收到写请求，更未必完成持久化 \\
应用语义 & 写了内存，不等于业务事务已经提交 \\
\bottomrule
\end{longtable}

这就是很多 mmap 程序出错的地方：它们看到另一个进程能读到新字节，就以为数据已经“保存到磁盘”。实际上那只是 page cache 可见性，不是掉电后的持久性。若随后机器掉电，而 dirty page 还没写回，磁盘上的文件可能仍是旧内容。若应用需要“写完记录后断电也能恢复”，至少要推进到 \code{msync(MS\_SYNC)} 或 \code{fsync}，并且还要考虑文件长度、目录项、日志或 manifest 的提交顺序。

\code{MAP\_PRIVATE} 的同一行代码完全不同。A 第一次写时，内核复制一页，把 A 的 PTE 指向匿名页；page cache 仍保存文件原始内容，B 也仍看到原始内容。此时“写入成功”只对 A 的进程地址空间成立，不对文件成立。私有映射不是“共享映射但晚点落盘”，而是写入后脱离文件页的匿名内存。

\subsection{回收与写回}

内存压力下，干净文件页可以直接回收，因为它们能从文件重新读入；脏文件页必须先写回；匿名页若没有 swap，就不能随意丢弃。

\begin{lstlisting}
reclaim_page(page):
  if page.refcount != 0 or page.locked:
    return BUSY
  if page.file_backed and not page.dirty:
    remove_from_page_cache(page)
    free(page)
    return OK
  if page.file_backed and page.dirty:
    enqueue_writeback(page)
    return RETRY_LATER
  if page.anon and swap_available:
    write_to_swap(page)
    free(page)
    return OK
  return BUSY
\end{lstlisting}

这个算法说明 page cache 不只是性能缓存，也是内存回收策略的一部分。错误地回收 dirty page 会丢数据；错误地保留所有文件页会让匿名内存被挤爆。

\section{测试}

\begin{enumerate}
\tightlist
\item 基本映射：映射文件第一页，读取字节应等于文件内容。
\item 懒加载：\code{mmap} 后不访问不应触发磁盘读，第一次访问才缺页。
\item 权限：只读 fd 不能建立可写 \code{MAP\_SHARED}。
\item 私有写：\code{MAP\_PRIVATE} 写入后文件内容不变。
\item 共享写：\code{MAP\_SHARED} 写入后 \code{read} 同区间能看到 page cache 更新。
\item \code{munmap}：解除映射后访问同地址触发 SIGSEGV 或等价错误。
\item 截断：文件被 truncate 后访问越界映射要触发 SIGBUS 或教学系统定义的文件映射错误。
\item 回写：脏页经过 \code{msync} 后，模拟块设备能观察到写请求。
\end{enumerate}

\begin{rubricbox}
本章合格标准：能画出 VMA、PTE、page cache、inode 的关系；能解释 private/shared 映射的写故障差异；能说明 \code{msync}、\code{fsync} 和业务提交不是同一个边界。
\end{rubricbox}

\subsection{本章压缩进来的持久化与恢复}

文件映射、page cache 和写回最终都要服务恢复语义。普通应用常把“别的进程现在能读到新字节”误认为“断电后仍能读到新字节”；存储系统则必须把每一步写成可恢复协议。主线里先抓住四层提交：内存可见、page cache 变脏、块设备完成、文件系统和目录项持久。业务提交通常还要在这些层之上，再写 manifest、日志或版本号。

\begin{longtable}{p{0.18\linewidth}p{0.42\linewidth}p{0.28\linewidth}}
\toprule
边界 & 已经保证什么 & 仍未保证什么 \\
\midrule
\code{memcpy} 到映射区 & 当前进程地址空间已改变 & 文件页未必写回，私有映射不改文件 \\
其他进程读到新字节 & 共享 page cache 可见 & 掉电后未必存在 \\
\code{msync(MS\_SYNC)} & 指定映射脏页写回并等待结果 & 目录项、rename、其他文件未必提交 \\
\code{fsync(file)} & 文件数据和必要元数据推进到持久层 & 多文件事务顺序未必正确 \\
\code{rename + fsync(dir)} & 目录项切换可恢复 & 应用协议仍要能识别旧版/新版 \\
\bottomrule
\end{longtable}

可靠更新通常需要单调证据。先写新数据并同步，再写临时 manifest 并同步，随后 \code{rename} 替换正式 manifest，最后同步目录。崩溃恢复时，程序只相信通过校验和、长度、版本号或事务记录证明完整的版本。若先更新 manifest 再同步数据，恢复程序可能看到“版本 42 已提交”，但数据页还停在旧内容。持久化错误最危险的地方是正常测试很难暴露，只有把崩溃点插在每个提交边界之后，才能证明恢复算法不会相信半成品。

\section{源码级补充：do\_mmap、XArray、writeback、O\_DIRECT 与 DAX}

\subsection{do\_mmap 到 mmap\_region}

Linux 的 \code{mmap} 路径会从系统调用参数进入 \code{do\_mmap}，经过权限、地址选择、长度和 flags 校验后，调用 \code{mmap\_region} 建立或合并 VMA。文件映射还会调用 \code{file->f\_op->mmap}，让具体文件系统或设备设置 \code{vm\_ops}。

\begin{lstlisting}
sys_mmap
  -> ksys_mmap_pgoff
  -> vm_mmap_pgoff
  -> do_mmap
  -> get_unmapped_area
  -> mmap_region
  -> file->f_op->mmap
\end{lstlisting}

读源码时看 \filepath{mm/mmap.c}、\filepath{mm/filemap.c}。重点关注 \code{vm\_flags}：\code{VM\_READ/WRITE/EXEC/SHARED/MAYWRITE/LOCKED/HUGEPAGE} 等标志共同决定 fault、mprotect、mlock、回收和写回行为。

\subsection{address\_space 与 XArray}

page cache 不挂在 fd 上，而挂在文件的 address\_space 上。一个 inode 的 \code{i\_mapping} 用 XArray 以页索引为 key 保存缓存页或 folio。这样不同进程、不同 fd、read/write/mmap 都能找到同一份文件页。

\begin{lstlisting}
filemap_fault(vmf):
  mapping = vmf.vma.file.f_mapping
  index = offset_to_page_index(vmf.address)
  folio = xa_load(mapping.i_pages, index)
  if folio missing:
    folio = page_cache_sync_readahead(...)
    read_folio_from_filesystem(...)
  install_pte(vmf.address, folio)
\end{lstlisting}

XArray 替代老的 radix tree，支持稀疏索引、并发查找和特殊 entry。读源码时不必先精通 XArray 实现，但要知道 page cache 的 key 是文件偏移，不是虚拟地址。

\subsection{readahead 与访问模式}

顺序读时，内核会预读后续页；随机读时，过度预读会浪费 I/O 和内存。应用可用 \code{madvise} 给出提示：\code{MADV\_SEQUENTIAL}、\code{MADV\_RANDOM}、\code{MADV\_DONTNEED}、\code{MADV\_HUGEPAGE}。

\begin{lstlisting}
on_page_cache_miss(file, index):
  if sequential_pattern_detected(file):
    submit_readahead(index + 1, window)
  else:
    read_single_page(index)
\end{lstlisting}

预读失败案例：数据库随机读大文件，错误地被当成顺序流，page cache 被无用页污染；或者媒体顺序读取未触发预读，设备队列深度上不去。性能报告要注明访问模式和缓存状态。

\subsection{writeback 与错误范围}

脏页写回由 flusher 线程、内存压力、\code{fsync/msync} 等触发。Linux 用 \code{writeback\_control} 描述写回目标，用 address\_space 记录写回错误。一个 \code{write} 成功后发生的异步写回错误，可能在之后的 \code{fsync} 才报告。

\begin{lstlisting}
fsync(file):
  filemap_fdatawrite(mapping)
  filemap_fdatawait(mapping)
  err = file_check_and_advance_wb_err(file)
  if err:
    return err
  return filesystem_fsync_metadata(file)
\end{lstlisting}

错误范围跟踪很重要。否则一个进程造成的写回错误可能被另一个不相关进程消费，或者永远没人知道。可靠程序必须检查 \code{fsync} 返回值。

\subsection{O\_DIRECT 和 DAX}

\code{O\_DIRECT} 尝试绕过 page cache，把用户页直接用于块 I/O。它通常要求地址、长度和文件偏移对齐设备或文件系统要求。绕过 page cache 不等于绕过一致性：若同一文件同时被 buffered I/O 和 direct I/O 访问，内核必须同步或失效相关缓存页。

DAX 面向持久内存，允许文件直接映射到持久介质地址，跳过 page cache。这样 load/store 可直接访问持久介质，但持久化顺序、cache flush 和失败原子性更需要应用和文件系统谨慎处理。

\begin{longtable}{p{0.2\linewidth}p{0.34\linewidth}p{0.34\linewidth}}
\toprule
模式 & 优点 & 风险 \\
\midrule
buffered I/O & page cache、预读、合并、易用 & 双重缓存、fsync 边界易误解 \\
\code{O\_DIRECT} & 减少缓存污染，数据库可控 & 对齐、短 I/O、混用一致性复杂 \\
DAX & 跳过 page cache，低延迟持久内存 & cache flush、持久顺序、硬件依赖强 \\
\bottomrule
\end{longtable}

\subsection{截断、SIGBUS 与映射失效}

文件映射后，文件可能被另一个进程 truncate。若映射区间超过新的文件大小，访问超出部分通常触发 \code{SIGBUS}，因为虚拟地址属于 VMA，但背后的文件页不再存在。这和访问完全未映射地址的 \code{SIGSEGV} 不同。

\begin{lstlisting}
file_truncate(inode, new_size):
  invalidate page cache pages beyond new_size
  unmap mappings beyond new_size or mark invalid
  wake waiters

fault_mapped_file(addr):
  if page_index beyond inode.size:
    signal(SIGBUS)
\end{lstlisting}

这个边界常被忽略。mmap 程序不能假设映射后文件大小永远不变，尤其是共享文件、日志文件和被管理进程替换的文件。

\subsection{msync、fsync 和 fdatasync}

\code{msync} 面向映射区间，把 dirty mapped pages 写回；\code{fsync} 面向文件，等待文件数据和必要元数据持久；\code{fdatasync} 通常只保证读取文件数据所需的元数据。它们重叠但不等价。

\begin{longtable}{p{0.2\linewidth}p{0.52\linewidth}}
\toprule
接口 & 语义重点 \\
\midrule
\code{msync} & 映射页写回，可按地址范围控制 \\
\code{fsync} & 文件数据和必要元数据持久化 \\
\code{fdatasync} & 文件数据和读取必需元数据，不一定同步所有时间戳 \\
\code{syncfs} & 某挂载文件系统的脏数据同步 \\
\bottomrule
\end{longtable}

数据库和存储服务必须明确使用哪个边界。只调用 \code{msync} 某段映射，不一定同步目录项或业务 manifest；只调用 \code{fsync} 文件，不一定说明应用的多文件事务已提交。

用一个索引文件更新看清楚持久化边界。程序维护两个文件：\code{data.dat} 保存记录，\code{manifest} 保存“当前有效数据文件和版本”。它用 mmap 修改 \code{data.dat} 的某一页，然后更新 manifest 指向新版本：

\begin{lstlisting}
memcpy(mapped_data + off, record, len);
msync(mapped_data + page_start, PAGE_SIZE, MS_SYNC);
write(manifest_fd, "version=42\n", 11);
fsync(manifest_fd);
\end{lstlisting}

这段看起来已经很谨慎，但还要继续问：manifest 文件是覆盖写还是 rename 新文件？目录项是否 fsync？data 文件长度是否已经扩展并持久？如果 \code{record} 跨两页，只 msync 了一页会怎样？如果 \code{msync} 成功但 \code{fsync(manifest)} 失败，恢复时应该相信哪个版本？持久化不是“调用一个同步函数”就完了，而是要把恢复算法需要的证据按顺序写到介质上。

\begin{longtable}{p{0.2\linewidth}p{0.46\linewidth}p{0.22\linewidth}}
\toprule
动作 & 它推进了哪层状态 & 仍未保证什么 \\
\midrule
写 mmap 内存 & 改了 page cache 或私有匿名页 & 未必到块设备 \\
\code{msync(MS\_SYNC)} 数据页 & 等待指定映射脏页写回 & 未必同步目录项或其他文件 \\
\code{fsync(data\_fd)} & 推进 data 文件数据和必要元数据 & 不代表 manifest 已提交 \\
\code{rename(tmp, manifest)} & 原子切换目录项视图 & 目录项持久化仍要看目录 fsync \\
\code{fsync(dirfd)} & 推进目录项更新 & 不修复上层协议顺序错误 \\
\bottomrule
\end{longtable}

可靠更新通常要让恢复路径有单调证据。例如先把新数据页写入并同步，再写临时 manifest，\code{fsync} 临时文件，\code{rename} 替换正式 manifest，最后 \code{fsync} 目录。崩溃后要么看到旧 manifest 指向旧数据，要么看到新 manifest 指向已经同步过的新数据。若顺序反过来，先提交 manifest 再写 data，恢复时可能读到“版本 42 已提交”但数据页仍是旧内容。这里的核心不是某个固定配方，而是每一步都要说明：崩溃发生在这一行之后，恢复程序会看到哪些对象状态。

\subsection{userfaultfd 和用户态缺页处理}

userfaultfd 允许用户态接管某些缺页事件，用于虚拟机迁移、按需加载和用户态分页。内核把 fault 事件交给用户态管理进程，后者填充页面并唤醒 faulting 线程。

\begin{lstlisting}
fault on registered range:
  kernel queues userfault event
  faulting thread sleeps
  pager reads event
  pager copies page data into address
  pager wakes faulting thread
\end{lstlisting}

这个机制说明 page fault 不只属于内核内部，也可以成为用户态系统软件的接口。但它要求严格权限和资源控制，否则恶意 pager 可以让线程永久睡眠或制造内存压力。
